\subsection{Controlling every congruence solution}

Integral cocycles need not lift every cocycle modulo a power of the
uniformizer. Completeness supplies the required uniform bound directly.

\begin{lemma}\label{shared-complete-solution-count}
Let $A$ be a complete discrete valuation ring, with uniformizer $t$ and
finite residue field $k$. Let $C$ be a complete separated torsion-free
$A$-module, let $C'$ be separated and torsion-free, and let
$\Phi:C\longrightarrow C'$ be linear. Put
$S_m=\Phi^{-1}(t^mC')$ and $K=\ker\Phi$. Suppose $S_1/tC$ is finite
and $B\subseteq K$ is a submodule with $K/B$ torsion. Then some $E$
satisfies
\[
 \#(S_m/t^mC)\leq
 \#\bigl((B+t^mC)/t^mC\bigr)\,|k|^E
 \qquad(m\geq0).
\]
\end{lemma}

\begin{proof}
The descending images of $S_m$ in $S_1/tC$ stabilize at some $m_0$.
Since $S_{r+1}\cap tC=tS_r$, induction lifts this stabilization to
every power of $t$: a member of $S_{m_0+j}$ can be approximated modulo
$t^jC$ by members of arbitrarily deep $S_m$. Choose successive
approximations whose differences lie in $t^{j+r}C$.
Completeness gives a limit in $K$, and cancellation in $C'$ gives
\[
                    S_{m_0+j}\subseteq K+t^jS_{m_0}.
\]
The kernel is closed and saturated, so it is complete and
$K\cap t^mC=t^mK$. Thus $K/tK$ is finite; complete Nakayama makes
$K$ finitely generated, and $K/B$ is finite. Reduction modulo $t$
and the same cancellation identity show inductively that every
$S_m/t^mC$ is finite. The displayed inclusion bounds its quotient by
the image of $B$, for $m=m_0+j$, by the fixed number
\[
             \#(K/B)\,\#(S_{m_0}/t^{m_0}C).
\]
The finitely many smaller $m$ are absorbed in $|k|^E$.
\end{proof}

For the deformation problem, $C$ is the module of functions whose
reductions are locally constant, and $\Phi$ lists the cocycle,
polarization and detecting equations. Finite-level cocycles have
function lifts; an integral cocycle lift is unnecessary. The residual
finiteness hypothesis and integral detection therefore bound all
modular solutions. Integral boundaries reduce into the full modular
boundary set, which only enlarges the right side.

The remaining detection statement concerns the paired first-order
representation over $k((t))$. Its detecting functional is the specified
coefficient of $\operatorname{tr}(z(\sigma)\rho_1(\sigma))$; the
first-order terms contain the independent odd extension classes.
Ordinary trace against the block-diagonal representation cannot replace
this functional. Fixing the point and detecting elements precedes the
constant $E$; the subsequent auxiliary level may depend on $m$.
